\chapter{Iteration 4 --- results}
\label{ch:iter4}

\section*{Task and inputs}

Same AAD task as Chapter~\ref{ch:results}: predict the attended speaker
from 30-s trials. Three experiments here, each varying \emph{only} the
EEG/gaze source or decoder architecture:

\begin{itemize}
  \item \textbf{4-class EEG decoder} (§\ref{sec:iter4-4class}): input =
        32-ch EEG at \SI{64}{\Hz}; supervision = 4 candidate attended
        envelopes (one per speaker). Target = pick the speaker. Gaze,
        IMU, and video are not used.
  \item \textbf{EEG$\leftrightarrow$gaze predictability}
        (§\ref{sec:iter4-pred}): input = 32-ch EEG (for EEG$\to$gaze) or
        the 20-D Tobii gaze summary (for gaze$\to$attended). Target =
        gaze time-courses, or the attended-speaker label.
  \item \textbf{Gaze-residualised EEG AAD}
        (§\ref{sec:iter4-resid}): input = 32-ch EEG \emph{after}
        channel-wise regression of a 15-feature Tobii gaze + IMU
        regressor set; same CCA-mel-28 decoder as Chapter~\ref{ch:results}
        otherwise. Target = attended-speaker hemisphere.
\end{itemize}

Per-subject 5-fold CV; pooled = unweighted mean across 16 subjects with
bootstrap 95\,\% CI. All 48 SLURM tasks completed on the
\code{PAS2966} CPU queue.

\section{4-class EEG speaker-identity decoder}
\label{sec:iter4-4class}

CCA between EEG lags ($-200\ldots+500$\,ms) and 28-band gammatone envelopes
of each of the 4 possible attended speakers (Device-1 L/R, Device-2 L/R).
Test trials are classified by \emph{argmax} canonical correlation across
the four envelope candidates. Per-subject 5-fold CV, 100 trials.

\subsection{Per-subject accuracy}

\begin{table}[ht]
\centering
\begin{tabular}{lrrr}
\toprule
Subject & 4-class & hemisphere & inner/outer \\
\midrule
1  & 0.240 & 0.490 & 0.480 \\
2  & 0.270 & 0.520 & 0.480 \\
3  & 0.260 & 0.430 & 0.610 \\
4  & 0.220 & 0.480 & 0.450 \\
5  & 0.310 & 0.560 & 0.520 \\
6  & 0.240 & 0.560 & 0.500 \\
7  & 0.330 & 0.500 & 0.610 \\
8  & 0.250 & 0.520 & 0.500 \\
9  & 0.250 & 0.490 & 0.540 \\
10 & 0.280 & 0.450 & 0.550 \\
11 & 0.250 & 0.500 & 0.550 \\
12 & 0.210 & 0.540 & 0.450 \\
13 & 0.230 & 0.450 & 0.510 \\
14 & 0.210 & 0.450 & 0.480 \\
15 & 0.263 & 0.505 & 0.545 \\
16 & 0.270 & 0.530 & 0.520 \\
\midrule
\textbf{Pooled} & \textbf{0.255}    & \textbf{0.498}    & \textbf{0.518} \\
95\,\% CI       & [0.240, 0.271] & [0.479, 0.517] & [0.496, 0.542] \\
Chance          & 0.250             & 0.500             & 0.500 \\
\bottomrule
\end{tabular}
\end{table}

All three task framings land at chance. The 4-class decoder does not
distinguish individual speakers; even the inner-outer binary barely
exceeds chance ($+1.8$\,pp). This is expected: within a device, the
stereo speech file is replicated on both channels with different
amplitudes --- the envelope shape is near-identical, so envelope
reconstruction cannot pick between the two channels of the same device.

Window sweep $\{1,2,4,8,16,30\}$\,s shows no gradient beyond
$\pm 0.01$, confirming the backbone is not merely data-starved.

\section{Cross-modal predictability (EEG $\leftrightarrow$ gaze)}
\label{sec:iter4-pred}

\subsection{EEG $\rightarrow$ gaze (ridge regression, held-out Pearson $r$)}

\begin{table}[ht]
\centering
\begin{tabular}{lrr}
\toprule
Target & mean $r$ & std \\
\midrule
gaze2d\_x (scene-projected horizontal) & \textbf{0.028} & 0.020 \\
gaze2d\_y (scene-projected vertical)   & 0.016 & 0.015 \\
gaze3d\_x (listener-ref horizontal, mm) & 0.009 & 0.011 \\
gaze3d\_y (vertical, mm)                & 0.000 & 0.012 \\
gaze3d\_z (depth, mm)                   & 0.002 & 0.008 \\
\bottomrule
\end{tabular}
\end{table}

EEG carries almost no linear predictive information about instantaneous
gaze: held-out Pearson $r$ peaks at $0.028$ for the horizontal
scene-projected gaze and collapses to zero in 3-D body-referenced space.

\subsection{CCA EEG $\leftrightarrow$ gaze (first-component correlation)}

\begin{table}[ht]
\centering
\begin{tabular}{lr}
\toprule
 & Mean first-component $r$ \\
\midrule
In-sample (training residuals) & 0.147 $\pm$ 0.092 \\
Held-out                        & \textbf{0.081} $\pm$ 0.063 \\
\bottomrule
\end{tabular}
\end{table}

A small but consistent shared subspace exists (held-out $r$ $= 0.081$),
but it is an order of magnitude below what each modality exposes about
the attended-speaker label on its own.

\subsection{Gaze $\rightarrow$ attended (logistic classifier)}

\begin{table}[ht]
\centering
\begin{tabular}{lrrr}
\toprule
Task              & Chance & Pooled & 95\,\% CI \\
\midrule
Hemisphere (L vs R) & 0.500 & \textbf{0.769} & [0.684, 0.851] \\
Inner vs outer      & 0.500 & 0.680          & [0.603, 0.762] \\
4-class identity    & 0.250 & 0.551          & [0.439, 0.668] \\
\bottomrule
\end{tabular}
\end{table}

Reproduces the corrected-mapping numbers from Chapter~\ref{ch:results}.

\section{Gaze-residualised EEG AAD --- the headline control}
\label{sec:iter4-resid}

For every trial, a 15-feature gaze/IMU regressor set
(\code{gaze2d\_x/y}, \code{gaze3d\_x/y/z}, per-eye direction, pupil
diameters, gaze velocities, \mbox{$|\text{accel}|{-}g$}, $|\text{gyro}|$)
is channel-wise ridge-regressed out of the EEG ($\lambda {=} 10^{-3}$)
\emph{before} the CCA mel-28 decoder sees it. Compares against the same
decoder on un-residualised EEG (\emph{baseline}). Per-subject 5-fold CV,
hemisphere task, full 30-s trial.

\begin{table}[ht]
\centering
\begin{tabular}{lrr}
\toprule
Subject & Baseline & Gaze-residualised \\
\midrule
1  & 0.560 & 0.580 \\
2  & 0.570 & 0.610 \\
3  & 0.380 & 0.460 \\
4  & 0.490 & 0.520 \\
5  & 0.510 & 0.570 \\
6  & 0.510 & 0.580 \\
7  & 0.580 & 0.580 \\
8  & 0.490 & 0.570 \\
9  & 0.560 & 0.560 \\
10 & 0.520 & 0.430 \\
11 & 0.510 & 0.530 \\
12 & 0.440 & 0.520 \\
13 & 0.510 & 0.590 \\
14 & 0.520 & 0.450 \\
15 & 0.424 & 0.505 \\
16 & 0.470 & 0.440 \\
\midrule
\textbf{Pooled}    & \textbf{0.503}    & \textbf{0.531} \\
95\,\% CI          & [0.477, 0.528]    & [0.503, 0.558] \\
\bottomrule
\end{tabular}
\end{table}

\noindent
Mean paired $\Delta = +0.028$\,pp. Wilcoxon signed-rank
$p = 0.084$, $n = 16$. 11 of 16 subjects improve after removing the
gaze-correlated variance.

\textbf{Headline result}: removing gaze-correlated variance from EEG
\emph{improves} AAD accuracy, not reduces it. This directly refutes the
``eye-tracking-in-disguise'' concern: EEG carries AAD-relevant signal
that is independent of gaze, and the gaze-linked component is
effectively noise relative to the TRF task.

\section{Extended fusion (gaze + 4-class-EEG-probabilities)}

Fusion inputs: Tobii gaze summary features (20-D) $+$ the 4-class CCA
softmax probabilities ($\text{p}_1, \ldots, \text{p}_4$). Classifier =
LightGBM ($n=200$ estimators) with 5-fold stratified CV per subject.
Three task framings; S15's 4-class EEG cache was missing so fusion is
computed on 15 subjects; gaze-only is pooled on all 16.

\begin{table}[ht]
\centering
\begin{tabular}{lrrr}
\toprule
Task              & Gaze-only & EEG-only (4-class probs) & Early fusion \\
\midrule
Hemisphere    & \textbf{0.770} & 0.501 & 0.749 \\
Inner vs outer & 0.690          & 0.477 & 0.643 \\
4-class identity & 0.547        & 0.255 & 0.493 \\
\bottomrule
\end{tabular}
\end{table}

\begin{itemize}
  \item Fusion does not exceed gaze-only on any task --- the 4-class EEG
        probabilities are too close to chance to add discriminative
        signal on top of a strong gaze feature vector.
  \item The \emph{only} bounded gain from fusion will come from either
        (a) a stronger EEG decoder (iter-5: deep models) or (b) subjects
        whose gaze is unreliable, where the EEG residual is additive.
  \item Per-subject fusion outperforms gaze-only only for 4 subjects
        on hemisphere task (S5: $+0.02$, S6: $+0.07$, S7: $+0.02$, S12:
        $+0.00$). These are the low-gaze-accuracy subjects --- exactly
        the regime where EEG is expected to help.
\end{itemize}

\section{What iter-4 settles}

\begin{enumerate}
  \item \textbf{Linear EEG alone cannot discriminate within-device L/R}
        (4-class $= 0.255$, chance $= 0.250$). The paradigm's stereo
        device design means envelope reconstruction is blind to channel
        identity. Iter-5 will test whether a deep decoder on
        \emph{raw EEG} (not reconstruction) can exploit the spatial
        lateralisation cues that CCA misses.
  \item \textbf{EEG has very low linear predictability of gaze}
        ($r \leq 0.03$) and a modest shared subspace
        ($r_\text{CCA, out} = 0.081$). Gaze contamination of EEG exists
        but is small.
  \item \textbf{Removing gaze \emph{improves} EEG-AAD by $+2.8$\,pp} ---
        the reviewer-ready answer to the ``eye-tracking in disguise''
        question.
  \item \textbf{The fusion story} needs a better EEG backbone to beat
        the gaze ceiling. The wiring is in place; iter-5 just swaps the
        EEG module.
\end{enumerate}
